\exercisesection{2}

\begin{enumerate}
\def\labelenumi{\arabic{enumi}.}
\tightlist
\item
  \begin{enumerate}
  \def\labelenumii{(\alph{enumii})}
  \tightlist
  \item
    Show that, if \(\mathfrak{a}\) is an ideal \(\neq 0\) in a Dedekind domain \(\mathbf{A}\) , the ring \(\mathbf{A} / \mathfrak{a}\) is a principal ideal ring (Algebra, Chapter VII, § 1, Exercise 5) (note that \(\mathbf{A} / \mathfrak{a}\) is semi-local and argue as in Proposition 1 of no. 2). Deduce again that every fractional ideal of \(\mathbf{A}\) is generated by at most two elements (cf.~Exercise 11 (a)).
  \end{enumerate}
\end{enumerate}

\begin{enumerate}
\def\labelenumi{(\alph{enumi})}
\setcounter{enumi}{1}
\tightlist
\item
  In the polynomial ring \(\mathbf{A} = k[\mathbf{X},\mathbf{Y}]\) over a field \(k\) , let \(\mathfrak{m}\) be the ideal \(\mathbf{AX} + \mathbf{AY}\) . Show that for every integer \(n\) the minimum number of generators of the ideal \(\mathfrak{m}^n\) is \(n + 1\) .
\end{enumerate}

\begin{enumerate}
\def\labelenumi{\arabic{enumi}.}
\setcounter{enumi}{1}
\tightlist
\item
  \begin{enumerate}
  \def\labelenumii{(\alph{enumii})}
  \tightlist
  \item
    Let \(\mathfrak{a}\) , \(\mathfrak{b}\) be two integral ideals of a Dedekind domain \(\mathbf{A}\) ; show that there exists an integral ideal \(\mathfrak{c}\) which is prime to \(\mathfrak{b}\) and such that \(\mathfrak{ac}\) is principal (use Proposition 9 of §1, no. 5).
  \end{enumerate}
\end{enumerate}

\begin{enumerate}
\def\labelenumi{(\alph{enumi})}
\setcounter{enumi}{1}
\tightlist
\item
  Let \(\mathfrak{a}, \mathfrak{b}\) be two integral ideals of \(\mathbf{A}\) ; show that there exists \(x \neq 0\) in the field of fractions of \(\mathbf{A}\) such that \(x\mathfrak{a}\) is an integral ideal of \(\mathbf{A}\) which is prime to \(\mathfrak{b}\) (same method). Deduce that the module \(\mathfrak{a}/\mathfrak{ab}\) is isomorphic to \(\mathbf{A}/\mathfrak{b}\) .
\end{enumerate}

\begin{enumerate}
\def\labelenumi{\arabic{enumi}.}
\setcounter{enumi}{2}
\item
  Let A be a Dedekind domain, K its field of fractions, P the set of prime ideals \(\neq 0\) of A and for all \(p \in P\) let \(v_{p}\) be the corresponding essential valuation on K. For every sub-A-module M of K and all \(p \in P\) we write \(v_{p}(M) = \inf v_{p}(x)\) (taken in \(\overline{\mathbf{R}}\) ); if \(M \neq 0\) , \(v_{p}(M) < +\infty\) for all \(p \in P\) and \(v_{p}(M) \leqslant 0\) for almost all \(p \in P\) . If M, N are two sub-A-modules of K, show that the relation \(M \subset N\) is equivalent to \(v_{p}(N) \leqslant v_{p}(M)\) for all \(p \in P\) (use Proposition 9 of §1, no. 5). Conversely, for every family \((v_{p})_{p \in P}\) of elements equal to an integer or \(-\infty\) and such that \(v_{p} \leqslant 0\) for almost all \(p \in P\) , there exists a unique sub-A-module M of K such that \(v_{p}(M) = v_{p}\) for all \(p \in P\) .
\item
  Let A be a Dedekind domain and K its field of fractions. If L is a subfield of K such that A is integral over \(A \cap L\) , then \(A \cap L\) is a Dedekind domain (show first that \(A \cap L\) is a Krull domain and then that every prime ideal of \(A \cap L\) is maximal, using Chapter V, §2, no. 1, Theorem 1).
\item
  \begin{enumerate}
  \def\labelenumii{(\alph{enumii})}
  \tightlist
  \item
    Let k be a field, K the field \(k(\mathbf{X}, \mathbf{Y})\) of rational functions in two indeterminates over k and A the polynomial ring K{[}Z{]} in one indeterminate, which is a principal ideal domain. Let L be the subfield \(k(\mathbf{Z}, \mathbf{X} + \mathbf{Y}\mathbf{Z})\) of \(k(\mathbf{X}, \mathbf{Y}, \mathbf{Z})\) ; show that \(A \cap L\) is not a Dedekind domain (prove that there are prime ideals which are not maximal and \(\neq 0\) in this ring).
  \end{enumerate}
\end{enumerate}

\begin{enumerate}
\def\labelenumi{(\alph{enumi})}
\setcounter{enumi}{1}
\tightlist
\item
  Let A be a Dedekind domain which is not a principal ideal domain and for which the ideal class group C(A) is finite * (the ring of integers of a finite algebraic extension of Q has the latter property). * Let \((\mathfrak{a}_{j})_{1 \leqslant j \leqslant r}\) be a representative system of the group I(A) of fractional ideals \(\neq 0\) modulo the subgroup of fractional principal ideals, consisting of integral ideals of A and let \(p_{i} (1 \leqslant i \leqslant s)\) be the prime ideals dividing at least one of the \(a_{j}\) . Let S (resp. T)
\end{enumerate}

denote the multiplicative set of elements \(x \in \mathbf{A}\) such that \(v_{\mathfrak{p}_i}(x) = 0\) for all \(i\) (resp. consisting of 1 and the \(x \in \mathbf{A}\) such that \(v_{\mathfrak{q}}(x) = 0\) for all the prime ideals \(\mathfrak{q}\) distinct from the \(\mathfrak{p}_i\) and \(v_{\mathfrak{p}_i}(x) \geqslant 1\) for all \(i\) ; the hypothesis implies that \(\mathbf{T}\) is not reduced to 1). Show that \(S^{-1}\mathbf{A}\) and \(T^{-1}\mathbf{A}\) are principal ideal domains and that \(A = (S^{-1}A) \cap (T^{-1}A)\) .

\begin{enumerate}
\def\labelenumi{\arabic{enumi}.}
\setcounter{enumi}{5}
\item
  Show that in a Dedekind domain the notions of primary ideal, irreducible ideal, primal ideal (Chapter IV, § 2, Exercise 33), quasi-prime ideal (Chapter IV, § 2, Exercise 34) and power of a prime ideal are identical.
\item
  Let A be a Noetherian domain. Show that the following properties are equivalent:
\end{enumerate}

\((\alpha)\) A is a Dedekind domain.

(β) For every maximal ideal m of A, there exists no ideal a distinct from m and \(m^{2}\) such that \(m^{2} \subset a \subset m\) .

\((\gamma)\) For every maximal ideal \(\mathfrak{m}\) of \(\mathbf{A}\) , the set of primary ideals for \(\mathfrak{m}\) is totally ordered by inclusion.

(δ) For every maximal ideal m of A, every primary ideal for m is a product of prime ideals.

(Prove that each of the properties \((\gamma)\) and \((\delta)\) implies \((\beta)\) ; then show that \((\beta)\) implies that \(A_{m}\) is a field or a discrete valuation ring (cf.~Chapter VI, § 3, no. 5, Proposition 9).)

Give an example of a local integral domain satisfying properties \((\beta)\) , \((\gamma)\) and \((\delta)\) , which is not a valuation ring (cf.~Chapter VI, § 3, Exercise 7).

¶ 8. Let A be an integral domain in which every prime ideal ≠0 is invertible. Show that A is a Dedekind domain. (Prove first that every prime ideal p' ≠ 0 of A is maximal, noting that, if p is a prime ideal such that p ≠ p' and p ⊃ p', then p': p = p' (Chapter II, § 1, Exercises 8 (b)) and on the other hand that p': p = p'p⁻¹, whence a contradiction is deduced. Deduce that A is Noetherian (Chapter II, § 1, Exercise 6 (b)) and finally apply Chapter II, § 5, no. 6, Theorem 4 to show that A\_m is a field or a discrete valuation ring for every maximal ideal m of A.)

\begin{enumerate}
\def\labelenumi{\arabic{enumi}.}
\setcounter{enumi}{8}
\tightlist
\item
  Let A be an integral domain.
\end{enumerate}

\begin{enumerate}
\def\labelenumi{(\alph{enumi})}
\item
  Let \((\mathfrak{p}_i)_{1\leqslant i\leqslant m},(\mathfrak{p}_j')_{1\leqslant j\leqslant n}\) be two finite families of invertible prime ideals such that \(\mathfrak{p}_1\mathfrak{p}_2\dots \mathfrak{p}_m = \mathfrak{p}_1'\mathfrak{p}_2'\dots \mathfrak{p}_n'\) . Show that \(m = n\) and that there exists a permutation \(\pi\) of \([1,n]\) such that \(\mathfrak{p}_i' = \mathfrak{p}_{\pi (i)}\) for all \(i\) such that \(1\leqslant i\leqslant n\) .
\item
  Let \(\mathfrak{p}\) be a prime ideal of \(\mathbf{A}\) and \(a\) an element of \(\mathbf{A} - \mathfrak{p}\) ; if \(\mathfrak{p} \subset \mathbf{A}a + \mathfrak{p}^2\) , show that \(\mathfrak{p} = \mathfrak{p}(\mathbf{A}a + \mathfrak{p})\) ; deduce that, if \(\mathfrak{p}\) is invertible, then \(\mathbf{A}a + \mathfrak{p} = \mathbf{A}\) .
\item
  Show that, if every ideal of A is a product (not necessarily a priori unique) of prime ideals of A, A is a Dedekind domain. (Show first that every invertible prime ideal p is maximal, considering an element \(a \in A - p\) and decompositions of \(Aa + p\) and \(Aa^{2} + p\) as products of prime ideals and applying (a)
\end{enumerate}

in the ring A/p to show that \(Aa^{2} + p = (Aa + p)^{2}\) ; then use (b). Prove then that every prime ideal \(p \neq 0\) is invertible, by decomposing as a product of prime ideals a principal ideal Ab where \(b \in p\) , observing that the factors of this product are invertible and, by applying the above, show that p is necessarily equal to one of these factors. Conclude with the aid of Exercise 8.)

¶ 10. Let A be a ring in which every ideal is a product of prime ideals.

\begin{enumerate}
\def\labelenumi{(\alph{enumi})}
\item
  Show that, for every prime ideal p of A, A/p is a Dedekind domain (Exercise 9).
\item
  Show that the set of minimal prime ideals \(\mathfrak{p}_i\) of \(\mathbf{A}\) is finite (decompose (0) as a product of prime ideals).
\item
  Suppose that \(\mathbf{A} / \mathfrak{p}_i\) is not a field. Show that, if \(y \in \mathfrak{p}_i\) , then, for all \(x \in \mathbf{A} - \mathfrak{p}_i\) , there exists \(z \in \mathfrak{p}_i\) such that \(y = zx\) . (Consider in A the decompositions as products of prime ideals of Ax and Ax + Ay and the corresponding decompositions in \(\mathbf{A} / \mathfrak{p}_i\) .) Deduce that, for all non-zero \(y \in \mathfrak{p}_i\) , \(Ay = \mathfrak{p}_i\) (consider \(\mathfrak{p}_i / Ay\) , decomposing Ay as a product of prime ideals); show finally that \(\mathfrak{p}_i = \mathfrak{p}_i^2\) and therefore that there exists in \(\mathfrak{p}_i\) an idempotent \(e_i\) such that \(\mathfrak{p}_i = Ae_i\) (cf.~Chapter II, § 4, Exercise 15). Then A is the direct composition of the ring \(Ae_i\) and the ring \(\mathbf{A}(1 - e_i)\) , isomorphic to the Dedekind domain \(\mathbf{A} / \mathfrak{p}_i\) .
\item
  Suppose now that all the \(p_{i}\) are maximal, so that A is a direct composition of rings of the form \(A/p_{i}^{v_{i}}\) (Chapter II, § 1, no. 2, Proposition 5) and we may therefore confine our attention to the case where A is primary or where the unique prime ideal p of A is nilpotent. Show that in this case the hypothesis implies that A is a principal ideal ring (Algebra, Chapter VII, § 1, Exercises 5 and 6).
\item
  Conclude from (c) and (d) that A is the product of a finite number of Dedekind domains and a principal ideal ring.
\end{enumerate}

\begin{enumerate}
\def\labelenumi{\arabic{enumi}.}
\setcounter{enumi}{10}
\tightlist
\item
  Let K be a field and \((v_{\iota})_{\iota\in I}\) a family of discrete valuations on K satisfying conditions \((\mathbf{AK}_{\mathbf{I}})\) and \((\mathbf{AK}_{\mathbf{III}})\) of §1, no. 3 and such that, for every integer r, every family \((n_{h})_{1\leqslant h\leqslant r}\) of integers \(\geqslant0\) , every family \((\iota_{h})_{1\leqslant h\leqslant r}\) of distinct elements of I and every family \((a_{h})_{1\leqslant h\leqslant r}\) of elements of K, there exists \(x\in K\) such that \(v_{\iota_{h}}(x-a_{h})\geqslant n_{h}\) for \(1\leqslant h\leqslant r\) and \(v_{\iota}(x)=0\) for ι distinct from the \(\iota_{h}\) . Show that the intersection A of the valuation rings of the \(v_{\iota}\) is a Dedekind domain whose field of fractions is K and that \((v_{\iota})_{\iota\in I}\) is the family of essential valuations of A. (Use Exercise 7 of §1; then prove that two distinct prime ideals of height 1 are relatively prime and deduce that every prime ideal \(\neq0\) of A is maximal.)
\end{enumerate}

¶ 12. Let A be an integral domain and K its field of fractions. Show that the following properties are equivalent:

\((\alpha)\) For every prime ideal \(\mathfrak{p}\) of A, the local ring \(A_{\mathfrak{p}}\) is a valuation ring.

(β) For every maximal ideal m of A, the local ring \(A_{m}\) is a valuation ring.

(γ) Every finitely generated ideal \(\neq0\) in A is invertible (and therefore the finitely generated fractional ideals \(\neq0\) form a group).

(δ) Every torsion-free finitely generated A-module is projective.

(ε) For all \(x \in \mathbf{K}\) , the fractional ideal \(\mathbf{A} + \mathbf{A}x\) is invertible.

(ζ) For all \(x \neq 0\) in \(\mathbf{K}\) , there exists \(y \in \mathbf{K}\) such that \(y \equiv 0 \pmod{1}\) , \(y \equiv 0 \pmod{x}\) and \(y \equiv 1 \pmod{1 - x}\) .

\((\eta)\) If \((a_{i})\) is a finite family of ideals of A, for the system of congruences \(x \equiv c_{i} \pmod{a_{i}}\) to have a solution it is necessary and sufficient that \(c_{i} \equiv c_{j} \pmod{(a_{i} + a_{j})}\) for every pair of indices i, j (``Chinese remainder theorem'').

(θ) If \(a, b, c\) are three ideals of \(A\) , then \(a \cap (b + c) = a \cap b + a \cap c\) .

\begin{enumerate}
\def\labelenumi{(\roman{enumi})}
\item[(ι)]
  If \(a, b, c\) are three ideals of \(A\) , then \(a + (b \cap c) = (a + b) \cap (a + c)\) .
\item[(κ)]
  A is integrally closed and every finitely generated ideal of A is divisorial.
\end{enumerate}

(λ) A is integrally closed and for \(z \neq 0\) in \(\mathbf{K}\) there exist \(x, y\) in \(\mathbf{A}\) such that \(z = x + yz^2\) .

(μ) A is integrally closed and, if a, b, c are three finitely generated fractional ideals \(\neq0\) , the relation ab = ac implies b = c.

\begin{enumerate}
\def\labelenumi{(\alph{enumi})}
\setcounter{enumi}{21}
\tightlist
\item
  For every finitely generated A-module M, the torsion submodule of M is a direct factor of M.
\end{enumerate}

(σ) Every ring B such that A ⊂ B ⊂ K is an integrally closed domain.

Then A is called a Prüferian domain (or a Prüfer domain).

(To prove the equivalence of \((\alpha)\) , \((\beta)\) , \((\gamma)\) and \((\delta)\) , use Chapter II, § 5, no. 2, Theorem 1 and no. 6, Theorem 4. To prove that \((\varepsilon)\) implies \((\gamma)\) , use the identity \((a + b)(b + c)(c + a) = (a + b + c)(ab + bc + ca)\) between three fractional ideals in an integral domain. To prove that \((\zeta)\) implies \((\eta)\) , argue by induction on the number of \(a_i\) . The equivalence of \((\eta)\) , \((\theta)\) and \((\iota)\) has been shown in Algebra, Chapter VI, § 1, Exercise 25. Prove that \((\lambda)\) implies \((\varepsilon)\) by noting that \((\lambda)\) implies the relation

\[
x (\mathbf {A} + \mathbf {A} z) \subset z (\mathbf {A} + \mathbf {A} z),
\]

using Exercise 6 (b) of § 1 and noting that

\[
(\mathbf {A} + \mathbf {A} z) (\mathbf {A} y + \mathbf {A} x / z) = \mathbf {A}.
\]

To prove that \((\mu)\) implies \((\lambda)\) noting that always

\[
z (\mathrm{A} + \mathrm{Az}) \subset (\mathrm{A} + \mathrm{Az} ^ {2}) (\mathrm{A} + \mathrm{Az}).
\]

Prove that \((\kappa)\) implies \((\lambda)\) noting that every fractional principal ideal \(A t\) which contained \(A\) and \(A z^2\) also contains \(A z\) . To prove that \((\nu)\) implies \((\gamma)\) , consider, for a finitely generated integral ideal \(\mathfrak{a} \neq 0\) of \(A\) , a non-zero element \(c \in \mathfrak{a}(A: \mathfrak{a})\) and apply Exercise 32 of §1 to the ideal \(\mathfrak{b} = c \mathfrak{a}\) . To prove that \((\sigma)\) implies \((\beta)\) , reduce it to the case where \(A\) is a local ring, consider an element \(z \in K - A\) and show that \(z^{-1} \in A\) , noting that the hypothesis implies that \(z \in A[z^2]\) .

¶ 13. (a) In a Prüferian domain A, let \(a\) , \(b\) be two finitely generated ideals; show that, if \(b \notin a\) , there exists a finitely generated ideal \(c\) such that \(a \subset c\) , \(a \neq c\) and \(bc \subset a\) (consider the ideal \(a + b\) ).

\begin{enumerate}
\def\labelenumi{(\alph{enumi})}
\setcounter{enumi}{1}
\item
  Deduce from (a) that, if a is a finitely generated ideal in a Prüferian domain A, then in the ring A/a every element which is not a divisor of zero is invertible. In particular, a prime ideal of A cannot be finitely generated if it is maximal.
\item
  Show that in a Prüferian domain A two prime ideals p, p' neither of which is contained in the other are relatively prime (consider, for every maximal ideal m of A, the ideal \(pA_{m} + p'A_{m}\) ).
\item
  In a Prüferian domain A let \(\mathfrak{a}\) be a finitely generated ideal. For \(\mathfrak{a}\) to be primal (Chapter IV, § 2, Exercise 33), it is necessary and sufficient that \(\mathfrak{a}\) be contained only in a single maximal ideal of A (use (b)); \(\mathfrak{a}\) is then irreducible and quasi-prime (Chapter IV, § 2, Exercise 34), so that, for finitely generated ideals, these three notions coincide. For \(\mathfrak{a}\) to be primary, it is necessary and sufficient that it be contained only in a single prime ideal (necessarily maximal) \(\mathfrak{m}\) ; for \(\mathfrak{a}\) to be strongly primary (Chapter IV, § 2, Exercise 27), it is necessary and sufficient also that \(\mathfrak{mA}_{\mathfrak{m}}\) be principal.
\end{enumerate}

\begin{enumerate}
\def\labelenumi{\arabic{enumi}.}
\setcounter{enumi}{13}
\item
  Let A be a Prüferian domain. Show that for an A-module M to be flat, it is necessary and sufficient that it be torsion-free (use Exercise 12 (δ)). Deduce that A is a coherent ring (Chapter I, § 2, Exercise 12).
\item
  \begin{enumerate}
  \def\labelenumii{(\alph{enumii})}
  \tightlist
  \item
    If A is a Prüferian domain, A/p is Prüferian for every prime ideal p of A.
  \end{enumerate}
\end{enumerate}

\begin{enumerate}
\def\labelenumi{(\alph{enumi})}
\setcounter{enumi}{1}
\item
  If A is Prüferian, so is \(S^{-1}A\) for every multiplicative subset S of A not containing 0.
\item
  Let K be a field and \((\mathbf{A}_{\lambda})\) a non-empty right directed family of subrings of K. Show that, if the \(A_{\lambda}\) are Prüferian domains, so is their union A (use Exercise 12 (ζ)).
\end{enumerate}

\begin{enumerate}
\def\labelenumi{\arabic{enumi}.}
\setcounter{enumi}{15}
\tightlist
\item
  Let A be a Prüferian domain, K its field of fractions and L a (finite or otherwise) algebraic extension of K; show that the integral closure A' of A in L is a Prüferian domain. (Let \(x \in L\) and let \(a_{0}X^{n} + a_{1}X^{n-1} + \cdots + a_{n}\) be a polynomial in A{[}X{]} of which x is a root; if we write
\end{enumerate}

\[
a _ {0} \mathrm{X} ^ {n} + a _ {1} \mathrm{X} ^ {n - 1} + \dots + a _ {n} = (\mathrm{X} - x) (b _ {0} \mathrm{X} ^ {n - 1} + \dots + b _ {n - 1}),
\]

show that the \(b_{i}\) and the \(b_{i}x\) belong to \(A'\) , using §1, Exercise 12 (a). If \(a = \sum_{i=0}^{n} A a_{i}\) , \(b = \sum_{j=0}^{n-1} B b_{j}\) , show then that the ideal \(b.B a^{-1}\) is the inverse of \(B + B x\) .)

¶ 17. Let A be an integral domain.

\begin{enumerate}
\def\labelenumi{(\alph{enumi})}
\item
  For A to be Bezoutian (§ 1, Exercise 20), it is necessary and sufficient that it be Prüferian and pseudo-Bezoutian (§ 1, Exercise 21).
\item
  For A to be a Dedekind domain, it is necessary and sufficient that it be a Prüferian Krull domain (show that every divisorial ideal of A is finitely generated, using Exercise 12 (κ) and § 1, Exercise 11 (a); deduce that every prime ideal \(\neq 0\) of A is maximal (Exercise 13 (b)), then that every prime ideal of A is finitely generated and conclude with the aid of Chapter II, § 1, Exercise 6).
\item
  For A to be a Dedekind domain, it is necessary and sufficient that it be Prüferian and strongly Laskerian (Chapter IV, § 2, Exercise 28). (Observe that, if A is Prüferian and strongly Laskerian, every maximal ideal m of A is such that \(A_{m}\) is a discrete valuation ring (Chapter VI, § 3, Exercise 8), then that for all non-zero \(x \in A\) there is a product of a finite number of maximal ideals of A contained in the ideal Ax and conclude that A is a Krull domain.)
\end{enumerate}

\begin{enumerate}
\def\labelenumi{\arabic{enumi}.}
\setcounter{enumi}{17}
\tightlist
\item
  Let A be a ring, the union of a right directed family of subrings \(A_{\alpha}\) which are Dedekind domains. Suppose that, for every prime ideal p of \(A_{\alpha}\) , there exists \(\beta \geqslant \alpha\) such that there are at least two distinct prime ideals of \(A_{\beta}\) lying above p.
\end{enumerate}

\begin{enumerate}
\def\labelenumi{(\alph{enumi})}
\item
  Show that the ring of all algebraic integers (the integral closure of \(\mathbf{Z}\) in \(\mathbf{C}\) ) satisfies the above conditions (cf.~Chapter V, § 2, Exercise 6).
\item
  Let \(v\) be a non-improper valuation on \(\mathbf{A}\) , \(z \in \mathbf{A}\) such that \(v(z) > 0\) and let \(\mathfrak{a}\) be the ideal of \(\mathbf{A}\) consisting of the \(x\) satisfying \(v(x) \geqslant v(z)\) . Show that \(\mathbf{A}: \mathfrak{a} = \mathbf{A}\) and hence that \(\mathfrak{a}\) is not invertible, although for every maximal ideal \(m\) of \(\mathbf{A}\) , \(\mathfrak{a}\mathrm{A}_{m}\) is a principal ideal in the valuation ring \(\mathrm{A}_{\mathfrak{m}}\) (cf.~Chapter II, § 5, no. 6, Theorem 4).
\item
  Deduce from (a) and (b) an example of a Prüferian domain which is completely integrally closed but not a Krull domain (cf.~Exercises 15 (c) and 17 (b) and Chapter V, § 1, Exercise 14).
\end{enumerate}

¶ 19. An integral domain A is called pseudo-Prüferian if the set of finitely generated divisors of D(A) (§ 1, Exercise 11) is a group. A pseudo-Prüferian domain is regularly integrally closed (§ 1, Exercise 29). A pseudo-Bezoutian domain (§ 1, Exercise 21) is pseudo-Prüferian; a Prüferian domain is pseudo-Prüferian; a Krull domain is pseudo-Prüferian.

\begin{enumerate}
\def\labelenumi{(\alph{enumi})}
\item
  Show by examples that the converses of the last three assertions are not necessarily true.
\item
  Let \(\theta\) be an irrational number \(>0\) and \(\Gamma\) the group \(\mathbf{R}^2\) ordered by taking the set of ordered pairs \((\alpha, \beta)\) such that \(\alpha \geqslant 0\) and \(\beta \geqslant \theta \alpha\) as the set of positive elements; the ordered group \(\Gamma\) is a complete lattice. Let B be the pseudo-principal domain derived from \(\Gamma\) by the procedure of § 1, Exercise 22 and let A be the subring the intersection of B and the algebra over \(k\) of the subgroup \(\mathbf{Q}^2\) of \(\mathbf{R}^2\) . Show that A is a completely integrally closed domain but is not pseudo-Prüferian (note that, if K is the field of fractions of A and U the group of invertible elements of A, the ordered group \(\mathrm{K}^*/\mathrm{U}\) is isomorphic to the group \(\Gamma \cap \mathbf{Q}^2\) , ordered by the ordering induced by that on \(\Gamma\) ).
\item
  If the domain A is pseudo-Prüferian, show that the polynomial domain A{[}X{]} is pseudo-Prüferian (cf.~§ 1, Exercise 30 (c)).
\item
  Let A be a pseudo-Prüferian domain and K its field of fractions. Show that the integral closure B of A in an algebraic extension L of K is a pseudo-Prüferian domain (method of Exercise 16, using § 1, Exercise 29 (a)).
\end{enumerate}

* (e) Deduce from Exercise 30 (d) of § 1 an example of an increasing sequence \((\mathbf{A}_n)\) of factorial domains with the same field of fractions whose union is not a regularly integrally closed domain (nor a fortiori a pseudo-Prüferian domain) (in the example quoted, take B to be a discrete valuation ring). *

\begin{enumerate}
\def\labelenumi{\arabic{enumi}.}
\setcounter{enumi}{19}
\tightlist
\item
  Let A be a Dedekind domain, K its field of fractions, L a finite algebraic extension of K and B the integral closure of A in L, which is a Dedekind domain. Let f be an ideal of B. For there to exist a ring C such that \(A \subset C \subset B\) and such that f is the conductor of B in A (Chapter V, § 1, no. 5), it is necessary and sufficient that, for every prime ideal \(p'\) of B containing f such that the field \(B/p'\) is isomorphic to \(A/(p' \cap A)\) (prime ideals of residue degree 1), the intersections of f and the transporter f: \(p'\) of \(p'\) in f with A be equal. (Note that, if there exists such a ring C, the conductor of B in the ring \(C_{0} = A + f\) is also equal to f; the existence of C is therefore equivalent to the fact that \(A + f\) contains no ideal of B distinct from f and containing f.~To prove that the latter condition is equivalent to that of the statement, reduce it to the case where A is a discrete valuation ring.)
\end{enumerate}

¶ 21. (a) Let A be an integral domain, f, g two polynomials in A{[}X{]} and h = fg. Let a, b, c denote the ideals of A generated respectively by the coefficients of f, g, h. Show that, if \(\deg(g) = n\) , then \(a^{n+1}b = a^{n}c\) . (Let \(f(X) = \sum_{i=1}^{m} a_{i}X^{i}\) , \(g(X) = \sum_{j=1}^{n} b_{j}X^{j}\) ; for every increasing sequence

\[
\sigma = (i _ {k}) _ {1 \leqslant k \leqslant n + 1}
\]

of integers \(\leqslant m\) and all \(j\) , let

\[
u _ {\sigma , j} = a _ {i _ {1}} a _ {i _ {2}} \dots a _ {i _ {j + 1}} b _ {j} a _ {i _ {j + 2}} \dots a _ {i _ {n + 1}};
\]

take on the set of \(u_{\sigma,j}\) a total ordering such that, for \(j < j'\) , \(u_{\sigma,j} < u_{\tau,j'}\) for all \(\sigma, \tau\) , and, for all \(j\) , \(u_{\sigma,j} \leqslant u_{\tau,j}\) if and only if \(\sigma \leqslant \tau\) in the lexicographical ordering on \([0, m]^{n+1}\) . Then argue by induction on this totally ordered set.)

\begin{enumerate}
\def\labelenumi{(\alph{enumi})}
\setcounter{enumi}{1}
\item
  Deduce from (a) that, if A is Prüferian (Exercise 12), then c = ab.
\item
  Taking A to be the polynomial ring \(\mathbf{Z}[\mathbf{Y}]\) , give an example of two polynomials \(f, g\) of the first degree in \(\mathrm{A}[\mathbf{X}]\) such that \(c \neq ab\) .
\end{enumerate}

¶ 22. (a) Let A be a Noetherian domain each of whose prime ideals ≠0 is maximal, so that every ideal a ≠ 0 may be written uniquely as a product \(\prod_{i} q_{i}\) of primary ideals relative to the distinct prime ideals \(p_{i}\) containing a.

Let \(\mathfrak{a}\) be an invertible ideal in \(\mathbf{A}\) and let \(\mathfrak{b}\) be any ideal \(\neq 0\) in \(\mathbf{A}\) ; show that there exists an ideal \(\mathfrak{c}\) such that \(\mathfrak{b} + \mathfrak{c} = \mathbf{A}\) and \(\mathfrak{ac}\) is principal (Observe that, if \((\mathfrak{p}_i)_{1 \leqslant i \leqslant n}\) is a finite family of distinct maximal ideals of \(\mathbf{A}\) and we write \(\mathfrak{r}_i = \prod_{j \neq i} \mathfrak{p}_j\) , then \(\mathfrak{a} = \sum_i \mathfrak{ar}_i\) and \(\bigcap_i \mathfrak{ar}_i = \mathfrak{ap}_1\mathfrak{p}_2 \ldots \mathfrak{p}_n\) ; using Chapter II, § 1, no. 2, Proposition 6, deduce the existence of an element \(x \in \mathfrak{a}\) such that \(x^{-1}\mathfrak{a} + \mathfrak{b} = \mathbf{A}\) .) Deduce that \(\mathfrak{a}\) is generated by two elements (cf.~Exercise 1).

\begin{enumerate}
\def\labelenumi{(\alph{enumi})}
\setcounter{enumi}{1}
\tightlist
\item
  Let K be a field and A the subring of the ring of formal power series \(K[[T]]\) consisting of the series \(a_{0} + T^{n}P(T)\) , where \(a_{0} \in K\) , \(P(T) \in K[[T]]\) , for a given integer n.~Show that A is a Noetherian local integral domain with a single prime ideal \(m \neq 0\) , but that the smallest cardinal of a system of generators of m is n.
\end{enumerate}

